\section{Introduction}
\label{sec:intro}

Open-source software is maintained by people who are free to leave at any time. Most contributors join for a short period, while a small core of developers carries most of the knowledge and the workload~\cite{joblin2017classifying,avelino2016novel}. When one of them leaves, issues remain unanswered, knowledge is lost, and the project itself may be abandoned~\cite{avelino2019abandonment}. For this reason, developer turnover has received considerable attention: researchers studied when developers take breaks and whether they come back~\cite{iaffaldano2019why,calefato2022will}, which activity patterns precede turnover~\cite{lin2017developer}, and why contributors give up~\cite{miller2019people}. In parallel, burnout research has linked exhaustion and disengagement of software professionals to job demands, organizational culture, and attrition~\cite{tulili2023burnout,trinkenreich2023model,trinkenreich2024attrition}.

While this body of knowledge provides compelling evidence of the relevance of turnover, a noticeable research gap remains. On the one hand, repository-based studies describe \textit{who} leaves through activity, tenure, and structural position, but they say little about the \textit{psychological} side of leaving, \ie how developers feel and relate to their community in the months before they disengage. On the other hand, the large body of work on developers' affect~\cite{novielli2019sentiment,calefato2018sentiment,imran2024uncovering} mines emotions from communication as a snapshot, without following them up to the moment developers leave~\cite{mantyla2016mining}. Moreover, most studies equate leaving a project with disengagement, although a developer who stops contributing to a project may simply move to another one. \textit{To the best of our knowledge, software engineering research still lacks evidence on whether psychological signals observable in developers' communication anticipate disengagement beyond structural and workload factors, and on whether they distinguish developers who leave a project from those who leave open source}.

We address this gap through the lens of the Job Demands-Resources (JD-R) model~\cite{demerouti2001job,bakker2007job}, which explains disengagement as the result of demands that are not balanced by resources. We focus on three signals that developers leave in their own words and in the reactions of their peers. First, \textit{emotions}~\cite{shaver1987emotion}, \eg anger and sadness, which reflect strain. Second, \textit{belonging}, \ie the identification with the community, which psycholinguistics associates with the use of first-person plural rather than singular pronouns, \ie \textit{we} rather than \textit{I}~\cite{tausczik2010psychological}. Third, \textit{recognition}, \ie the gratitude and the positive reactions that developers receive for their work, which the JD-R model treats as a resource~\cite{bakker2007job}.

\steObjectiveBox{\faBullseye\ \textbf{Goal.} In this paper, we investigate whether \textbf{psychological signals} in developers' communication, \ie emotions, belonging, and received recognition, \textbf{anticipate disengagement} from open-source projects beyond structural and workload factors, and whether they distinguish leaving a project from leaving open source.}

To this aim, we built a mining pipeline that reconstructs, for each developer and month, the established structural and workload measures used in turnover research, \ie core-peripheral role~\cite{joblin2017classifying}, truck-factor membership~\cite{avelino2016novel}, and off-hours activity~\cite{claes2018programmers}, together with the three psychological signals. Since existing emotion classifiers do not transfer well across platforms~\cite{novielli2020can,imran2022data}, we fine-tuned a RoBERTa model~\cite{liu2019roberta} on public gold standards, which outperforms existing tools on a public GitHub benchmark (micro-F1 0.62 \textit{vs.} 0.42). Moreover, for every developer who leaves a project, we checked whether they remained active elsewhere on GitHub. We then analyzed \tbd{N} developers of \tbd{N} open-source projects through event-time trajectories and survival models with time-varying covariates.

\tbd{The results show that [placeholder: main finding of RQ\textsubscript{1}]. [Placeholder: main finding of RQ\textsubscript{2}, \eg which signals remain associated with disengagement once role and workload are controlled for]. Finally, [placeholder: main finding of RQ\textsubscript{3}]. In a preliminary sample of 100 projects, 72\% of the developers who abandoned a project were still active elsewhere on GitHub in the following year.}

To sum up, our paper offers the following contributions: (i) an operationalization of emotions, belonging, and recognition as longitudinal, developer-level signals mined from repositories; (ii) an emotion classifier for GitHub communication that improves over existing tools on a public benchmark; (iii) empirical evidence on how these signals evolve before disengagement and on their association with it beyond structural and workload factors; (iv) a distinction between leaving a project and leaving open source, with evidence on how often the two differ; and (v) a publicly available replication package~\cite{replication} to facilitate future research and reproducibility.
